% Slide 1: introduction, rather than a separate title slide.
\begin{frame}[t]{Why compression matters}
  \begin{columns}[T,onlytextwidth]
    \column{0.76\textwidth}
      {\Large EEG Compression}\par\smallskip
      \small Debarghya Das \enspace/\enspace 240150049
    \column{0.19\textwidth}
      \raggedleft\includegraphics[width=\linewidth]{images/logo/IITG.jpg}
  \end{columns}
  \vspace{0.5cm}
  \begin{columns}[T,onlytextwidth]
    \column{0.40\textwidth}
      {\fontsize{28}{32}\selectfont\color{eegteal}+18.2\%}\par
      \small average annual growth in mobile broadband traffic, 2021--2026
    \column{0.54\textwidth}
      \large Video, speech and sensors\par\medskip
      \normalsize More data for LLMs, robotics and multimodal systems.
  \end{columns}
  \vspace{0.5cm}
  \begin{center}
    \begin{tikzpicture}[node distance=0.65cm]
      \node[stage,text width=2.25cm] (data) {Create data};
      \node[stage,text width=2.25cm,right=of data] (store) {Store it};
      \node[stage,text width=2.25cm,right=of store] (move) {Move it};
      \draw[link] (data)--(store); \draw[link] (store)--(move);
    \end{tikzpicture}
  \end{center}
  \small Compression reduces the bytes we must store and transmit.
  \credit{Source: \href{https://www.itu.int/itu-d/reports/statistics/wp-content/uploads/sites/5/2026/11/ITUFactsAndFigures2026.pdf}{ITU, Facts and Figures 2026, p. 16}. Traffic growth is a connectivity indicator, not total data creation.}
  \note{Compression matters because recording, storage and communication all consume resources. The ITU reports 18.2 percent average annual growth in mobile broadband traffic since 2021. Its 2026 estimate is 1.7 zettabytes. These figures describe network traffic, not the entire amount of data created worldwide. LLMs and robots give intuitive examples of systems that handle multiple data types. The question becomes more delicate when the data are brain recordings.}
\end{frame}

% Slide 2: all technique families from the website.
\begin{frame}[t]{Compression across modalities}
  \begin{columns}[T,onlytextwidth]
    \column{0.46\textwidth}
      \textbf{Recover every value}\par\medskip
      \textcolor{eegblue}{Huffman coding}\\
      \small Frequent symbols get shorter codes.\par\medskip
      \normalsize\textcolor{eegblue}{ZIP / DEFLATE}\\
      \small Repeated strings + Huffman coding.\par\medskip
      \normalsize\textcolor{eegblue}{PNG}\\
      \small Pixel prediction + lossless coding.
    \column{0.49\textwidth}
      \textbf{Allow some information loss}\par\medskip
      \textcolor{eegcoral}{MP3}\\
      \small Spend bits using a hearing model.\par\smallskip
      \normalsize\textcolor{eegcoral}{Uniform / $\mu$-law quantization}\\
      \small Fewer amplitude levels.\par\smallskip
      \normalsize\textcolor{eegcoral}{JPEG}\\
      \small Transform + quantize image detail.\par\smallskip
      \normalsize\textcolor{eegcoral}{Thresholding / dithering}\\
      \small Reduce pixel precision; retain tone differently.
  \end{columns}
  \vfill
  \small\alert{Quantization discards precision. Entropy coding packs what remains.}
  \credit{Background: \href{https://www.rfc-editor.org/rfc/rfc1951}{Deutsch, RFC 1951 (1996)}; project Techniques module. Thresholding and dithering are precision-reduction steps, not complete file formats.}
  \note{These are the technique families in the website. Huffman coding and ZIP are general-purpose lossless ideas; PNG applies lossless coding to images. DEFLATE combines LZ77 string matching with Huffman coding. MP3 is perceptual audio coding, not an EEG compressor. Mu-law companding is invertible before quantization: the loss comes from subsequently using finitely many levels. PNG is lossless; JPEG is normally lossy. Thresholding and dithering can map grayscale to one bit per pixel; bit packing and file coding are still needed to store it compactly.}
\end{frame}

% Slide 3: a licensed illustration plus a technical acquisition diagram.
\begin{frame}[t]{Neural signal capture: EEG}
  \begin{columns}[c,onlytextwidth]
    \column{0.35\textwidth}
      \centering\includegraphics[width=0.95\linewidth]{images/eeg-cap.png}
    \column{0.60\textwidth}
      \large Electrodes measure\par
      \alert{voltage differences}\par\medskip
      \normalsize One time series per channel.\par\medskip
      Population-level activity, mixed with noise and artifacts.
  \end{columns}
  \vspace{0.45cm}
  \begin{center}
    \begin{tikzpicture}[node distance=0.3cm]
      \node[stage,text width=2.45cm] (electrode) {Scalp\\electrodes};
      \node[stage,text width=2.45cm,right=of electrode] (adc) {Amplify + sample};
      \node[stage,text width=2.45cm,right=of adc] (samples) {Channel $\times$ time};
      \draw[link] (electrode)--(adc); \draw[link] (adc)--(samples);
    \end{tikzpicture}
  \end{center}
  \small Time and frequency views reveal different signal features.
  \credit{Acquisition context: \href{https://doi.org/10.1093/gigascience/giz002}{Lee et al. (2019)}. EEG-cap icon: \href{https://commons.wikimedia.org/wiki/File:EEG_Cap_-_The_Noun_Project.svg}{CIV / The Noun Project, CC0}. Illustration is not our apparatus.}
  \note{EEG is non-invasive scalp recording. Each channel measures voltage relative to another location or reference; it is not an isolated neuron's output or a direct reading of thoughts. Amplification, filtering and analog-to-digital conversion produce sampled multichannel time series. Scalp position and reference affect the measurements. Frequency content, changes across channels and responses aligned to stimuli provide different analysis features. The icon is a generic illustration, not a photograph of OpenBMI participants or equipment.}
\end{frame}

% Slide 4: explicit illustrative arithmetic, not a measured study result.
\begin{frame}[t]{Why EEG needs compression}
  \begin{center}
    \large Many channels $\times$ continuous recording
    \vspace{0.4cm}

    {\large $64\times256\times16\times86\,400\,/\,8$}\par
    \small channels\quad samples/s\quad bits/sample\quad seconds/day
    \vspace{0.35cm}

    {\fontsize{27}{31}\selectfont\color{eegteal}2.83 GB per day}
  \end{center}
  \vspace{0.35cm}
  \begin{columns}[T,onlytextwidth]
    \column{0.30\textwidth}\centering\textbf{Storage}\\\small Long recordings accumulate.
    \column{0.32\textwidth}\centering\textbf{Wireless links}\\\small Bandwidth is limited.
    \column{0.31\textwidth}\centering\textbf{Wearable devices}\\\small Sending bits costs energy.
  \end{columns}
  \credit{Illustrative raw 16-bit recording; decimal GB, metadata excluded. Context: \href{https://doi.org/10.1109/ATC.2015.7388309}{Dao et al. (2015)} and \href{https://doi.org/10.3390/info13040186}{Zhang et al. (2022)}. Our benchmark uses a different 12-channel, float32 representation.}
  \note{This is arithmetic, not a measured recording size or a reported compression ratio: 64 channels times 256 samples per second times 16 bits times 86,400 seconds divided by eight equals 2,831,155,200 bytes. Larger montages and longer recordings scale the burden. For wireless wearable EEG, communication can also constrain bandwidth and energy. We did not measure battery savings or device energy. Our benchmark stores twelve channels at 128 Hz in float32, so the later storage allowances use a separate, explicitly defined baseline.}
\end{frame}
